\chapter{Measles and Mumps Serology}

\section*{What Is This Test?}
Measles and mumps serology detects IgG and IgM antibodies against measles (rubeola) virus and mumps virus to determine immune status, diagnose acute infection, or confirm past infection or vaccination. Both are single-stranded negative-sense RNA viruses of the \textit{Paramyxoviridae} family. Measles is among the most contagious infectious diseases (R0 of 12--18), transmitted via respiratory droplets and airborne spread, with an incubation period of 7--14 days. Clinical measles features the classic triad of cough, coryza, and conjunctivitis (the "3 Cs"), Koplik spots (tiny white lesions on the buccal mucosa appearing 1--2 days before rash), followed by a maculopapular rash that begins on the face and spreads cephalocaudally. Complications include otitis media (7--10\%), pneumonia (1--6\%, the most common cause of death), encephalitis (1/1,000), and subacute sclerosing panencephalitis/SSPE (4--11/100,000 cases). Mumps presents with fever, malaise, and parotitis (swelling of the parotid glands, often bilateral), with complications including orchitis (20--30\% of postpubertal males), oophoritis (5\% of postpubertal females), meningitis (1--10\%), and encephalitis (0.1\%).

\section*{Alternative Names}
Measles Antibody Panel; Rubeola Serology; Mumps IgG/IgM; MMR Immunity Panel; Measles Immunity Testing; Parotitis Serology; Measles and Mumps Antibody Titers

\section*{Principle}
Enzyme-linked immunosorbent assay (ELISA) or chemiluminescent immunoassay (CLIA) detects serum IgM and IgG antibodies against measles and mumps virus antigens. Measles IgM appears 1--3 days after rash onset (rarely before) and remains detectable for 30--60 days. Measles IgG appears within 5--10 days of rash onset, peaks at 2--3 weeks, and persists for life. Mumps IgM is detectable within 5--7 days of symptom onset and remains positive for 4--6 weeks. Mumps IgG appears within 7--10 days, peaks at 2--4 weeks, and typically persists for life, though levels may decline. IgG avidity testing (measles) distinguishes primary infection (low avidity) from past infection or vaccination (high avidity). Measles RT-PCR from throat swab, nasopharyngeal swab, or urine is the preferred method for acute diagnosis (collected within 7 days of rash onset). Mumps RT-PCR from buccal/oral swab or CSF is the preferred acute diagnostic method (collected within 3 days of parotitis onset). Virus isolation in cell culture (Vero/SLAM cells) confirms both infections but is slow (1--2 weeks). Plaque reduction neutralization test (PRNT) is the gold standard for measles antibody measurement but is labor-intensive and performed only at reference laboratories.

\section*{History}
Measles was described by the Persian physician Rhazes in the 9th century. The measles virus was first isolated by John Enders and Thomas Peebles in 1954. The measles vaccine (live attenuated, Edmonston B strain) was licensed in 1963; the current Moraten strain has been in use since 1968. Mumps virus was identified in 1934, and the mumps vaccine (Jeryl Lynn strain) was licensed in 1967. The combined MMR vaccine was introduced in 1971, with the second dose recommended in 1989. Measles was declared eliminated from the United States in 2000, but outbreaks have recurred due to imported cases in undervaccinated communities (667 cases in 2014, 1,282 cases in 2019). IgG ELISA tests for measles and mumps immunity screening largely replaced clinical history and prior vaccination records in the 1990s--2000s for healthcare worker screening and pre-employment clearance. RT-PCR revolutionized outbreak investigation at the turn of the 21st century, enabling rapid case confirmation and molecular epidemiology.

\section*{Why This Test Is Ordered}
\begin{itemize}
  \item Assessment of measles and mumps immunity status in healthcare workers, college students, international travelers, and military recruits; documentation of positive IgG is accepted as evidence of immunity in place of vaccination records
  \item Prenatal screening for measles and mumps immunity (rubella is always included in the obstetric panel as part of MMR screening)
  \item Confirmation of suspected acute measles in a patient with febrile rash illness, cough, coryza, conjunctivitis, especially with known exposure, international travel, or lack of vaccination
  \item Diagnosis of acute mumps in patients with parotitis, especially in the context of outbreaks in close-contact settings (colleges, sports teams)
  \item Evaluation of suspected measles encephalitis (CSF measles IgG index and PCR) in patients with acute neurologic deterioration and recent measles exposure
  \item Evaluation of suspected mumps meningitis or orchitis when parotitis is absent or atypical
  \item Outbreak investigation: serology plus PCR for rapid confirmation of measles and mumps in suspected outbreak settings to guide public health interventions (ring vaccination, quarantine of susceptible contacts)
\end{itemize}

\section*{Is This a Primary or Secondary Test}
Serum measles/mumps IgM is a primary diagnostic test for suspected acute infection when measured 3--28 days after measles rash onset or 5 days to 6 weeks after mumps symptom onset. Measles/mumps IgG is a primary test for determination of immunity. However, for acute diagnosis, RT-PCR of respiratory specimens/throat swabs (measles) or buccal swabs (mumps) is the preferred primary test due to higher sensitivity and the ability to genotype the virus for molecular epidemiology.

\section*{How This Test Is Performed}
Venous blood is collected in a serum separator tube (SST, red-top or gold-top tube). For IgM and IgG testing, 5--7 mL of whole blood yields sufficient serum. The specimen should be centrifuged and serum separated within 2 hours of collection; serum can be refrigerated at 2--8\textdegree{}C for up to 7 days or frozen at -20\textdegree{}C for longer storage. For measles PCR, a throat or nasopharyngeal swab (Dacron/nylon flocked, placed in viral transport medium) and a urine specimen (10--50 mL, clean catch) should be collected within 7 days of rash onset. For mumps PCR, a buccal/oral swab (massaging the parotid duct area for 30 seconds before swabbing the duct orifice near the upper second molar) collected within 3 days of parotitis onset, coupled with a throat swab, provides optimal sensitivity. Molecular specimens should be transported refrigerated (2--8\textdegree{}C) and processed within 72 hours.

\section*{What Preparation Is Needed}
No special patient preparation required for blood draw or specimen collection. Patients should avoid eating or drinking 30 minutes before buccal swab collection for mumps PCR.

\section*{Contraindications and Precautions}
Recent administration of MMR vaccine (within 6 weeks) may cause a positive measles or mumps IgM result (vaccine-induced IgM response), which should be interpreted cautiously. IgG may be positive within 5--14 days of MMR vaccination but does not necessarily indicate long-term immunity unless the post-vaccination titer is measured at least 4--6 weeks later. False-negative measles IgM occurs when the specimen is collected within the first 72 hours of rash onset (before IgM is detectable). False-positive measles IgM may occur in the presence of rheumatoid factor, other viral infections (parvovirus B19, EBV), or autoimmune disease. Mumps IgM has lower sensitivity than PCR, particularly in previously vaccinated individuals who may not mount an IgM response; a negative IgM does NOT exclude mumps infection in a vaccinated patient. IgG assays may have reduced sensitivity for detecting vaccine-induced immunity (compared to neutralizing antibody levels measured by PRNT) because the ELISA targets predominantly nucleoprotein rather than neutralizing antibodies directed against the hemagglutinin and fusion surface glycoproteins.

\section*{Risks}
Standard venipuncture risks: minor pain, bruising, hematoma, vasovagal syncope. There are no risks from the serologic or PCR tests themselves. The major public health risk is failure to diagnose acute measles, as one infected individual can transmit to 12--18 susceptible contacts in a non-immune population.

\section*{What Happens After the Test}
IgG/IgM serology results available in 1--24 hours. PCR results available in 4--24 hours. Positive measles or mumps IgM must be confirmed by RT-PCR where available, and the case must be reported to the local public health department immediately (both are nationally notifiable diseases). The patient with suspected measles should be placed in airborne precautions (negative-pressure room, N95 respirator) and public health investigation initiated to identify and potentially vaccinate susceptible contacts. For mumps, droplet precautions should be implemented for 5 days after onset of parotitis. Susceptible exposed contacts (those without evidence of immunity) should be offered MMR vaccine (within 72 hours of exposure for measles, which may provide some protection) or immune globulin (within 6 days for measles).

\section*{Understanding Results}

\subsection*{Below Normal}
\begin{itemize}
  \item \textbf{Measles IgG negative (nonreactive):} Absence of detectable measles IgG antibody indicates susceptibility to measles (no evidence of prior infection or adequate vaccination). The individual should receive MMR vaccine (2 doses separated by at least 28 days if no contraindications) unless contraindicated. A negative IgG in a patient with suspected acute measles does NOT exclude the diagnosis if measured before day 5 of the rash (IgG hasn't yet developed)
  \item \textbf{Mumps IgG negative (nonreactive):} Indicates susceptibility to mumps. MMR vaccine is recommended. Due to lower sensitivity of some commercial mumps IgG assays compared to neutralization assays, low-positive or borderline results may indicate waning immunity rather than true susceptibility
  \item \textbf{Measles IgM negative:} Timing is critical. An IgM collected 72 hours before rash onset through day 3 of rash may be falsely negative (30--50\% false-negative rate in the first 72 hours). If clinical suspicion remains, repeat IgM after 3--5 days or perform RT-PCR
  \item \textbf{Mumps IgM negative:} Does not exclude acute mumps, particularly in previously vaccinated individuals (up to 50\% of PCR-confirmed mumps in vaccinated individuals have negative IgM). RT-PCR of buccal swab is preferred for acute diagnosis
\end{itemize}

\subsection*{Above Normal}
\begin{itemize}
  \item \textbf{Measles IgM positive (reactive):} Indicative of acute measles infection when collected >3 days after rash onset in a patient with compatible clinical illness (fever, cough, coryza, conjunctivitis, maculopapular rash). Should be confirmed by RT-PCR. False positives may occur (parvovirus B19, EBV, rheumatoid factor) and confirmation is recommended
  \item \textbf{Measles IgG positive (reactive):} Indicates immunity from past infection or vaccination. IgG generally persists for life. A positive IgG does not distinguish between infection- and vaccine-induced immunity. For individuals with documented MMR vaccination, positive measles IgG is sufficient evidence of immunity regardless of the quantitative titer
  \item \textbf{Mumps IgM positive (reactive):} Suggestive of acute mumps infection. Confirm with RT-PCR. In vaccinated populations, mumps IgM sensitivity is moderate (50--70\%), and false positives may occur
  \item \textbf{Mumps IgG positive (reactive):} Indicates past infection or vaccination. IgG may wane over time; a low positive result in a person at high risk of exposure (outbreak setting) may warrant a booster MMR dose
  \item \textbf{Measles RT-PCR positive:} Confirms acute measles infection. Genotyping the N-450 region of the nucleoprotein gene identifies the viral genotype (A through D, with 24 recognized genotypes) for molecular epidemiologic tracking. Most US outbreaks are genotype D8 or B3
  \item \textbf{Mumps RT-PCR positive (buccal swab, CSF):} Confirms acute mumps infection. Genotyping of the SH gene discriminates mumps genotypes (A through N) and identifies the Jeryl Lynn vaccine strain (genotype A) vs. wild-type strains (genotype G currently predominant in the US)
\end{itemize}

\section*{Normal Results}
Measles: \textbf{Measles IgG Positive} \textbf{(immune)}, \textbf{Measles IgM Negative} \textbf{(no acute infection)}. Mumps: \textbf{Mumps IgG Positive} \textbf{(immune)}, \textbf{Mumps IgM Negative} \textbf{(no acute infection)}. A positive IgG indicates immunity. Results are qualitative; quantitative titers may be reported as index values or IU/mL but a specific protective titer threshold is not standardized for commercial ELISAs (PRNT titers $\geq$1:120 are considered protective).

\section*{Follow-up Tests}
\begin{itemize}
  \item \textbf{Measles RT-PCR (throat swab, NP swab, urine):} Confirmatory testing for all suspected acute measles cases; collect within 7 days of rash onset; throat/NP swab has highest sensitivity
  \item \textbf{Mumps RT-PCR (buccal swab):} Confirmatory testing for all suspected acute mumps; collect within 3 days of parotitis onset; collect from the parotid duct orifice (Stensen duct) near the upper second molar
  \item \textbf{Measles IgG avidity testing:} Distinguishes primary infection (low avidity, indicating recent infection) from remote infection or vaccination (high avidity); performed at reference laboratories, useful in equivocal cases
  \item \textbf{Measles PRNT (plaque reduction neutralization test):} Gold standard for measuring functional neutralizing antibodies; used as a confirmatory test for immunity when commercial IgG ELISA results are equivocal (e.g., borderline positive) in outbreak settings; titer $\geq$1:120 is considered protective
  \item \textbf{CSF analysis (measles IgG index, PCR):} When measles encephalitis (ADEM or measles inclusion body encephalitis) is suspected; elevated CSF measles IgG index suggests intrathecal antibody synthesis
  \item \textbf{Molecular genotyping:} Sanger sequencing of the N-450 hypervariable region (measles) or SH gene (mumps); performed at CDC/state public health laboratories for outbreak linkage and to distinguish wild-type from vaccine-strain virus
\end{itemize}

\section*{References}
\begin{enumerate}
  \item CDC. Measles. In: Epidemiology and Prevention of Vaccine-Preventable Diseases (The Pink Book). 14th ed. CDC; 2021.
  \item Moss WJ. Measles. Lancet. 2017;390(10111):2490-2502.
  \item Hviid A, Rubin S, Mühlemann K. Mumps. Lancet. 2008;371(9616):932-944.
  \item CDC. Prevention of measles, rubella, congenital rubella syndrome, and mumps: ACIP recommendations. MMWR Recomm Rep. 2013;62(RR-04):1-34.
  \item Rota PA, Moss WJ, Takeda M, et al. Measles. Nat Rev Dis Primers. 2016;2:16049.
\end{enumerate}

\clearpage
